\section{Conclusion}\label{sec:conclusion}

The main unconditional bounds of
\Cref{thm:t2,thm:main-upper} bracket the divisibility antichain
saturation game as
\[
  c\,\frac{n(\log\log n)^2}{\log n}
  \le
  \L(n)
  \le
  \left(\frac{\Wfour}{2}+o(1)\right)n.
\]
for an absolute constant $c>0$.  The lower bound comes from a weighted
activation strategy and an exact-move potential on residual rank-three fibers,
while the upper bound improves the elementary linear picture by using a
lower-profile envelope for Shortener's legal-prime prefix and a fourth-order
Bonferroni comparison.  The independent fan strategy gives the more explicit
but smaller lower bound
$(1/8-o(1))n\log\log n/\log n$.

The old safe-edge potential still fails in a fiber-supported auxiliary state,
but the failing reply is an artificially unscored exact-edge deletion.  In the
actual game that reply contributes one move, exactly balancing the lost
live-edge coefficient.  This distinction upgrades the rank-three construction
from conditional to unconditional; it does not resolve the linear-versus-
sublinear dichotomy.

On the structural side, the shield reduction and Theorem A show that short
upper-half shield prefixes cannot prove a linear game lower bound through the
weighted shadow functional $\beta_n(P)$.

Three obstructions in \Cref{sec:obstructions} explain why the next improvements
are unlikely to come from purely local relaxations, fixed shadow-cover
dichotomies, or separator-only strategies.  A sharper upper bound will need to
use dynamic information about the actual divisibility game, not merely the
static incidence geometry of low-rank shadows.

The central question remains the sharp order of growth of $\L(n)$, and in
particular whether $\L(n)=o(n)$.  The methods here give new constants, exact
auxiliary structure, and a set of formal targets for future work, but they do
not settle that asymptotic dichotomy.
